\subsection{A shorter useful life can change two different histories}
\label{sec:error-estimate-bridge}

A carrying amount does not reveal whether the estimate behind it was sound.
Suppose an administrative asset costs 120, becomes available for use on
1 January 2023 and has no residual value. The company records straight-line
depreciation of 10 in each of 2023 and 2024, leaving PPE of 100. On 1 April
2025, management concludes that the asset will be retired on 31 December
2028. That conclusion alone cannot tell us whether to change future expense
or correct the earlier accounts. We need the information that supported the
old estimate, including what was available when those accounts were
authorized.

Consider two explicitly synthetic histories. In the first, twelve years was
a reasonable useful life at acquisition and at both year-end reviews.
A newly enacted operating restriction on 1 April 2025 now limits use to
the end of 2028. That restriction neither existed nor was anticipated at the
earlier dates. In the second history, the reliable, approved schedule at
acquisition already specified six years, but a transcription error entered
twelve. The original schedule was available and could reasonably have been
obtained and considered at each authorization. Six years remains appropriate;
the new discovery exposes the old mistake. These histories differ in facts,
not in an accounting election available on the same known evidence.

AASB 108.5 makes that availability condition part of the definition of a
prior-period error. Paragraphs 32--34A distinguish new information and
developments from correction of an error; paragraphs 36--38 apply estimate
changes to the current and relevant future periods
\citep{aasb108comp2022}. AASB 116.50--51 and 60--62 require systematic
allocation, review of useful life and a method reflecting consumption
\citep{aasb116dated}. Here there is one component, even consumption, no
capitalization into another asset, and an independently satisfied
recoverability assessment. The shorter life does not replace that assessment.

Dates determine which accounts are being corrected. The 2024 statements
were authorized on 31 March 2025; the discovery occurs the following day.
The 2023 statements were authorized on 31 March 2024. No 2025 interim
statements have been issued. Thus the error affects two authorized annual
reports and an unissued current quarter. Moving discovery before the 2024
authorization would change the treatment of that report under paragraph 41.
The example concerns newly arising April conditions in the estimate history;
it makes no separate determination about events revealing conditions that
existed at the earlier reporting date.

The historical method comparison uses the retained December 2022 AASB 108
and 116 compilations and September 2023 AASB 112 compilation. It does not
establish a later amendment chain or actual-company applicability. Assume
a private, for-profit Tier 1 entity with no publicly traded ordinary or
potential ordinary shares, no filing for public issuance and no voluntary
EPS presentation. Those facts place this illustration outside the inspected
AASB 133.2--3 scope \citep{aasb133scope2021}. A public or voluntarily
presenting entity needs the EPS restatement as well; the financing example's
share schedule cannot simply be borrowed for this company.

\subsubsection{New information changes the remaining allocation}

Use equal monthly depreciation. In the estimate history the first three
months of 2025 still consume \(10\times3/12=2.5\), leaving 97.5 on
1 April. There are 45 months from April 2025 through December 2028. The
remaining carrying amount, rather than original cost or a reconstructed
earlier estimate, is allocated across those months:
\[
 d_{\mathrm{Apr-Dec\ 2025}}=\frac{97.5}{45}\times9=19.5,
 \qquad d_{2025}=2.5+19.5=22.
\]
Year-end PPE is \(100-22=78\). Each full year from 2026 through 2028
has depreciation \(97.5\times12/45=26\); the future allocations
\(19.5+26+26+26=97.5\) exhaust the April balance. The reported 2023
and 2024 depreciation remains 10 in each year. Revising those amounts to
20 would turn later information into a fact supposedly known earlier.

The estimate disclosure states the new restriction, the revised remaining
life and the affected amounts under AASB 108.39--40. Relative to continuing
the old schedule, 2025 depreciation increases by 12. Depreciation in each of
2026--2028 increases from 10 to 26; after retirement it is zero instead of
the old schedule's 10 in each of 2029--2034. These changes redistribute the
same total depreciable cost. They are consequences of the supplied retirement
date, not forecasts inferred from the single reported balance.

\subsubsection{An old error changes comparative profit and opening equity}

In the error history, correct annual depreciation has always been
\(120/6=20\). Assume the error is material and all period-specific and
cumulative effects are practicably determinable. The 2025 report presents
2024 as its comparative year. AASB 108.42 requires correction of that
comparative and of the opening balances for errors before its start;
paragraph 46 excludes the prior-period correction from 2025 profit.
We therefore need two linked amounts: the error accumulated before
1 January 2024 and the additional error during 2024.

Tax makes the equity effect smaller than the PPE correction. Supply a
conditional tax law that permits deductions of 20 annually in 2023--2028,
with an enacted rate of 25\% throughout the relevant recovery periods.
The asset is recovered through use, and initial book and tax bases are
both 120. There are no other differences or uncertain positions. Each year
the company earns cash operating income of 50 before depreciation; current
tax is paid on the following 31 March. These are authored tax assumptions,
not statements of Australian tax law. AASB 112.7, 15--17, 46--47 and 51
support the tax-base and rate method; paragraph 53 prohibits discounting
the deferred tax \citep{aasb112dated}.

The originally reported 2023 PPE of 110 exceeded tax base 100 by 10,
giving a deferred tax liability (DTL) of 2.5. In 2024, reported PPE 100
exceeded tax base 80 by 20, giving DTL 5. Correct PPE is instead 100 and
80 respectively, equal to the tax bases, so those DTL balances disappear.
Original pretax profit of 40 carried current tax 7.5 and deferred expense
2.5, leaving profit 30. Correct pretax profit of 30 carries current tax
7.5 and no deferred expense, leaving profit 22.5 in each year. Cash tax
does not change: the tax deductions were correct throughout.

At the earliest comparative opening, 1 January 2024, reduce PPE by 10,
DTL by 2.5 and retained earnings by 7.5. During comparative 2024, the extra
depreciation of 10 and removal of deferred expense 2.5 reduce that year's
profit by another 7.5. Thus the cumulative 1 January 2025 adjustment is
PPE minus 20, DTL minus 5 and retained earnings minus 15. The opening
and comparative movements sum to the same cumulative correction; they
are not two separate postings of it.

To see the balance-sheet consequence, start with capital and cash of 200,
then pay 120 for the asset at acquisition. The following figures retain
the original publication alongside its corrected comparative. Initial
capital remains 200; all later equity changes here are retained profit.

\begin{center}
\small
\begin{tabular}{lrrrr}
\toprule
& 2024 original & 2024 restated & 2025 error & 2025 estimate\\
\midrule
Depreciation & 10 & 20 & 20 & 22\\
Profit after tax & 30 & 22.5 & 22.5 & 21\\
Cash & 172.5 & 172.5 & 215 & 215\\
Net PPE & 100 & 80 & 60 & 78\\
Current tax payable & 7.5 & 7.5 & 7.5 & 7.5\\
DTL & 5 & 0 & 0 & 4.5\\
Total equity & 260 & 245 & 267.5 & 281\\
\bottomrule
\end{tabular}
\end{center}

The estimate column starts from the valid original equity of 260. Its
2025 pretax profit is \(50-22=28\). Tax base is 60, so PPE 78 gives
DTL 4.5, a decrease of 0.5 from the opening DTL. Current tax 7.5 less
that deferred benefit gives tax expense 7 and profit 21. Compared with
the old 2025 schedule, the estimate reduces pretax profit by 12 and
after-tax profit by 9. The error column instead starts from corrected
equity of 245 and adds current profit 22.5. Equal cash of 215 therefore
does not establish equal earnings or equal equity.

\subsubsection{The discovery entry joins the books to the restatement}

The system discovers the mistake in April; it does not rewrite the time at
which it learned the fact. Its correction records 1 April as the posting
and knowledge date while attributing the prior-period amount to opening
equity. Debit retained earnings 15, debit DTL 5 and credit accumulated
depreciation 20. AASB 112.61A(b) and 62A(a) place the associated tax
correction directly in equity with the prior error, rather than crediting
current tax expense \citep{aasb112dated}.

There is also a current-quarter error. Before discovery the books have
already charged depreciation 2.5 and deferred tax expense 0.625 for
January--March. Correct depreciation for those three months is 5. Debit
current depreciation expense 2.5 and credit accumulated depreciation 2.5;
debit DTL 0.625 and credit current deferred tax expense 0.625. That
second tax entry reverses the current quarter's erroneous DTL increase.
It has a different origin from the prior-period DTL correction of 5.
After both corrections, April PPE is 75 and DTL is zero. Depreciation
15 over the remaining nine months leaves PPE 60; full-year current
depreciation is \(2.5+2.5+15=20\), not 40.

The comparative is a presentation of the corrected earlier periods, not
an instruction to post the correction a second time. A correction identity
binds the entity, affected periods, original close, accepted treatment and
target restatement. An exact retry returns the original receipt before
testing later staleness. A changed payload under that identity is a
conflict; a new request identifier cannot post the same target correction
again. Original publications remain retrievable and explicitly identified
as containing the error. These are proposed semantic compositions of the
existing time, dependency and commit rules, not new portable request
encodings or evidence of a running accounting engine.

The cash bridge provides a useful independent check. Cash after acquisition
is 80; 2023 income increases it to 130. Subsequent annual income of 50,
less payment of the previous year's tax 7.5, gives 172.5 and 215. Under
the case's operating-cash convention, corrected 2025 profit 22.5 plus
noncash depreciation 20 gives 42.5. The estimate history gives the same
cash increase from profit 21 plus depreciation 22 less the deferred-tax
benefit 0.5. Neither the opening-equity correction nor the comparative
restatement creates cash.

\subsubsection{When the classification cannot yet be made}

Missing original-life documentation is not proof that management made an
error, and finding a later shorter-life estimate is not proof that the
old estimate was reasonable. The unresolved case retains both questions
until evidence distinguishes them. A change of measurement basis requires
the policy-change procedure. Paragraph 35's instruction for difficulty
distinguishing a policy change from an estimate change does not resolve
an error-versus-estimate ambiguity.

Nor does a missing local file establish impracticability. Paragraphs 43--45
distinguish inability to determine period-specific effects from inability
to determine the cumulative effect; paragraphs 50--53 require the relevant
effort and prohibit hindsight. Such a case needs its own supported starting
date and correction route. The fully practicable example above cannot
supply them. Unknown materiality also leaves the material-restatement
route unestablished; it does not authorize ignoring a known error.

The error disclosure identifies the transcription mistake, comparative
depreciation and tax/profit effects, the affected closing balance-sheet
items and the earliest opening adjustment under paragraph 49. Cash and
current tax are unchanged. If EPS is required, its effect must be added
using that entity's actual numerator and share history. Missing tax law,
original availability or accepted correction treatment blocks the related
net-equity or posting claim. Balanced entries alone cannot supply those facts.
